\documentclass[../DoAn.tex]{subfiles}
\begin{document}

Phần này trình bày đặc tả chi tiết cho các use case phụ của hệ thống Smart Grading không được đặc tả trong mục 2.3 của Chương 2 do hạn chế về dung lượng. Mỗi đặc tả gồm tên use case, luồng sự kiện chính, luồng sự kiện phát sinh, tiền điều kiện và hậu điều kiện, theo đúng mẫu đặc tả đã sử dụng trong Chương 2.

\section{Đặc tả use case Quản lý trường học}
\subsection*{Tên use case}
Quản lý trường học.

\subsection*{Tiền điều kiện}
Quản trị viên hệ thống đã đăng nhập với vai trò admin.

\subsection*{Luồng sự kiện chính}
\begin{enumerate}
    \item Quản trị viên truy cập chức năng "Quản lý trường học".
    \item Hệ thống hiển thị danh sách các trường hiện có cùng các thông tin: tên trường, mã trường, địa chỉ, số lượng giáo viên, số lượng học sinh.
    \item Quản trị viên chọn chức năng "Tạo trường mới".
    \item Hệ thống hiển thị form tạo trường với các trường: tên trường, mã trường, địa chỉ, số điện thoại, email, cấu hình thang điểm.
    \item Quản trị viên nhập thông tin và nhấn "Lưu".
    \item Hệ thống kiểm tra tính hợp lệ của thông tin, lưu trường mới vào cơ sở dữ liệu và hiển thị thông báo thành công.
    \item Quản trị viên có thể chọn một trường để xem chi tiết, cập nhật thông tin hoặc xóa trường.
\end{enumerate}

\subsection*{Luồng sự kiện phát sinh}
\begin{itemize}
    \item Tại bước 6, nếu mã trường đã tồn tại: hệ thống thông báo và yêu cầu nhập mã khác.
    \item Tại bước 7, nếu trường đang có giáo viên hoặc học sinh hoạt động: hệ thống cảnh báo và yêu cầu xác nhận trước khi xóa.
\end{itemize}

\subsection*{Hậu điều kiện}
Nếu thành công, trường mới được tạo và có thể được gán cho quản trị viên trường.

\section{Đặc tả use case Quản lý giáo viên}
\subsection*{Tên use case}
Quản lý giáo viên.

\subsection*{Tiền điều kiện}
Quản trị viên hệ thống hoặc quản trị viên trường đã đăng nhập.

\subsection*{Luồng sự kiện chính}
\begin{enumerate}
    \item Người dùng truy cập chức năng "Quản lý giáo viên".
    \item Hệ thống hiển thị danh sách giáo viên theo trường (nếu là school-admin) hoặc toàn hệ thống (nếu là admin).
    \item Người dùng chọn "Thêm giáo viên".
    \item Hệ thống hiển thị form thêm giáo viên với các trường: họ tên, email, mã giáo viên, môn phụ trách.
    \item Người dùng nhập thông tin và nhấn "Lưu".
    \item Hệ thống tạo tài khoản giáo viên, gửi email kích hoạt đến giáo viên.
    \item Người dùng có thể chọn một giáo viên để xem chi tiết, cập nhật thông tin, phân công môn hoặc vô hiệu hóa tài khoản.
\end{enumerate}

\subsection*{Luồng sự kiện phát sinh}
\begin{itemize}
    \item Tại bước 6, nếu email đã tồn tại: hệ thống thông báo và yêu cầu sử dụng email khác.
    \item Tại bước 7, nếu vô hiệu hóa tài khoản: hệ thống yêu cầu xác nhận vì giáo viên có thể đang phụ trách lớp đang hoạt động.
\end{itemize}

\subsection*{Hậu điều kiện}
Nếu thành công, tài khoản giáo viên được tạo và có thể đăng nhập vào hệ thống sau khi kích hoạt qua email.

\section{Đặc tả use case Quản lý học sinh}
\subsection*{Tên use case}
Quản lý học sinh.

\subsection*{Tiền điều kiện}
Quản trị viên trường hoặc giáo viên đã đăng nhập.

\subsection*{Luồng sự kiện chính}
\begin{enumerate}
    \item Người dùng truy cập chức năng "Quản lý học sinh".
    \item Hệ thống hiển thị danh sách học sinh theo lớp hoặc theo trường.
    \item Người dùng chọn "Thêm học sinh vào lớp".
    \item Hệ thống cho phép thêm một học sinh mới (tạo tài khoản) hoặc thêm học sinh đã có tài khoản vào lớp.
    \item Người dùng nhập thông tin cần thiết và nhấn "Lưu".
    \item Hệ thống lưu thông tin và gửi thông báo cho học sinh.
    \item Người dùng có thể chọn một học sinh để xem chi tiết, cập nhật thông tin, chuyển lớp hoặc xóa khỏi lớp.
\end{enumerate}

\subsection*{Luồng sự kiện phát sinh}
\begin{itemize}
    \item Tại bước 5, nếu mã số sinh viên đã tồn tại: hệ thống thông báo và đề xuất sử dụng học sinh hiện có.
    \item Tại bước 7, nếu chuyển lớp: hệ thống yêu cầu chọn lớp đích và xác nhận.
\end{itemize}

\subsection*{Hậu điều kiện}
Nếu thành công, học sinh được thêm vào lớp và có thể tham gia các bài kiểm tra của lớp.

\section{Đặc tả use case Quản lý bài thi của học sinh}
\subsection*{Tên use case}
Quản lý các bài thi của học sinh.

\subsection*{Tiền điều kiện}
Giáo viên đã đăng nhập, có bài làm đã được chấm tự động trên hệ thống.

\subsection*{Luồng sự kiện chính}
\begin{enumerate}
    \item Giáo viên truy cập chức năng "Bài nộp" và chọn bài kiểm tra tương ứng.
    \item Hệ thống hiển thị danh sách bài làm của học sinh với các thông tin: tên học sinh, mã số, điểm, trạng thái.
    \item Giáo viên chọn một bài làm để xem chi tiết.
    \item Hệ thống hiển thị thông tin chi tiết: ảnh phiếu đã chấm, đáp án của học sinh, đáp án đúng, điểm từng câu, cảnh báo (nếu có).
    \item Giáo viên có thể chỉnh sửa thủ công các câu bị nhận dạng sai bằng cách chọn lại đáp án đúng.
    \item Giáo viên nhấn "Lưu chỉnh sửa" để cập nhật điểm.
    \item Giáo viên có thể nhấn "Phê duyệt kết quả" để chuyển bài làm sang trạng thái sẵn sàng công bố.
\end{enumerate}

\subsection*{Luồng sự kiện phát sinh}
\begin{itemize}
    \item Tại bước 5, nếu giáo viên chỉnh sửa: hệ thống ghi nhận lịch sử chỉnh sửa và cập nhật điểm tự động.
    \item Tại bước 7, nếu bài làm chưa được phê duyệt: học sinh chưa thể xem điểm.
\end{itemize}

\subsection*{Hậu điều kiện}
Nếu thành công, bài làm được cập nhật và sẵn sàng công bố cho học sinh khi giáo viên phát hành kết quả.

\section{Đặc tả use case Xem bài thi và điểm số}
\subsection*{Tên use case}
Xem bài thi và điểm số.

\subsection*{Tiền điều kiện}
Học sinh đã đăng nhập, bài thi đã có kết quả được công bố.

\subsection*{Luồng sự kiện chính}
\begin{enumerate}
    \item Học sinh truy cập chức năng "Bài thi của tôi" hoặc "Điểm của tôi".
    \item Hệ thống hiển thị danh sách các bài thi đã tham gia với điểm tổng và trạng thái.
    \item Học sinh chọn một bài thi để xem chi tiết.
    \item Hệ thống hiển thị: điểm tổng, điểm từng câu, đáp án của học sinh, đáp án đúng, ảnh phiếu đã chấm.
    \item Học sinh có thể xem các đơn phúc khảo đã gửi cho bài thi này và trạng thái của chúng.
    \item Học sinh có thể tải về báo cáo điểm cá nhân dạng PDF.
\end{enumerate}

\subsection*{Luồng sự kiện phát sinh}
\begin{itemize}
    \item Tại bước 4, nếu bài thi có câu bị giáo viên chỉnh sửa thủ công: hệ thống hiển thị thông báo để học sinh hiểu rõ.
    \item Tại bước 6, nếu tải về thất bại: hệ thống thông báo và cho phép thử lại.
\end{itemize}

\subsection*{Hậu điều kiện}
Nếu thành công, học sinh nắm được kết quả chi tiết bài thi của mình.

\section{Đặc tả use case Tạo câu hỏi, đề thi}
\subsection*{Tên use case}
Tạo câu hỏi, đề thi.

\subsection*{Tiền điều kiện}
Giáo viên đã đăng nhập, có quyền quản lý ngân hàng câu hỏi.

\subsection*{Luồng sự kiện chính}
\begin{enumerate}
    \item Giáo viên truy cập chức năng "Ngân hàng câu hỏi".
    \item Hệ thống hiển thị danh sách câu hỏi hiện có theo môn học và chủ đề.
    \item Giáo viên chọn "Tạo câu hỏi mới".
    \item Giáo viên chọn một trong ba phương thức: tạo thủ công, sinh bằng AI, hoặc nhập từ file.
    \item Đối với tạo thủ công: hệ thống hiển thị form với các trường nội dung, loại câu hỏi, các đáp án, đáp án đúng, độ khó, chủ đề, giải thích.
    \item Đối với sinh bằng AI: hệ thống hiển thị form yêu cầu nhập chủ đề và số lượng câu hỏi; giáo viên nhấn "Sinh câu hỏi"; hệ thống gọi dịch vụ AI và hiển thị kết quả; giáo viên có thể chỉnh sửa trước khi lưu.
    \item Đối với nhập từ file: hệ thống cho phép tải lên file Excel/CSV theo mẫu.
    \item Giáo viên lưu câu hỏi vào ngân hàng; câu hỏi được đánh dấu là "chờ phê duyệt" nếu do AI sinh.
\end{enumerate}

\subsection*{Luồng sự kiện phát sinh}
\begin{itemize}
    \item Tại bước 6, nếu dịch vụ AI gặp lỗi: hệ thống thông báo và cho phép thử lại hoặc chuyển sang tạo thủ công.
    \item Tại bước 7, nếu file nhập không đúng định dạng: hệ thống thông báo lỗi cụ thể và cho phép tải mẫu.
\end{itemize}

\subsection*{Hậu điều kiện}
Nếu thành công, câu hỏi được lưu vào ngân hàng và có thể sử dụng để tạo đề thi.

\section{Đặc tả use case Đăng ký/Đăng nhập}
\subsection*{Tên use case}
Đăng ký/Đăng nhập.

\subsection*{Tiền điều kiện}
Người dùng truy cập được vào ứng dụng web hoặc di động.

\subsection*{Luồng sự kiện chính}
\begin{enumerate}
    \item Người dùng truy cập trang đăng nhập.
    \item Người dùng nhập email và mật khẩu, nhấn "Đăng nhập".
    \item Hệ thống kiểm tra thông tin đăng nhập.
    \item Nếu hợp lệ, hệ thống cấp access token và refresh token, chuyển hướng đến dashboard phù hợp với vai trò.
    \item Nếu người dùng chưa có tài khoản, có thể chọn "Đăng ký" và làm theo quy trình đăng ký gồm nhập thông tin cá nhân, xác thực email.
    \item Nếu người dùng quên mật khẩu, có thể chọn "Quên mật khẩu" và làm theo quy trình khôi phục qua email.
\end{enumerate}

\subsection*{Luồng sự kiện phát sinh}
\begin{itemize}
    \item Tại bước 3, nếu thông tin sai: hệ thống thông báo lỗi và cho phép thử lại.
    \item Tại bước 4, nếu tài khoản chưa xác thực email: hệ thống yêu cầu xác thực trước.
    \item Tại bước 6, nếu email không tồn tại trong hệ thống: hệ thống thông báo và không gửi email khôi phục.
\end{itemize}

\subsection*{Hậu điều kiện}
Nếu thành công, người dùng được đăng nhập và truy cập các chức năng theo vai trò.

\end{document}